
\section{Numerical example}
\label{sec:numerical-example}


To illustrate the methods presented, the 2-DOF system shown in figure
\ref{fig:duf_schematic} is used throughout the thesis. It will be referred to as
the coupled duffing system, with parameters listed in table \ref{tab:duf_par}
and natural frequencies and damping ratio for the two modes listed in table
\ref{tab:duf_eigen}. The equations of motion are

\begin{equation}
  \label{eq:2dof}
  \begin{aligned}
    &m_1\ddot x_1 + c \dot x_1 + kx_1 + \mu x_1^3 + d(x_1 - x_2) = p \\
    &m_2\ddot x_2 + c \dot x_2 + kx_2 + \mu x_2^3 + d(x_2 - x_1) = 0 \\
  \end{aligned}
\end{equation}
where $x_1,x_2, p$ all depends on time.

An example on how a real continous beam with different nonlinear boundary
conditions is transformed into a coupled MDOF nonlinear duffing oscillator is
given in \textcite{mhermansen2017a}


\begin{center}
  \begin{tabular}{*{6}{c}}
    \hline
    $m_1$ (kg) & $m_2$ (kg) & $k$ (N/m) & $d$ (N/m) & $c$ (N/ms) & $\mu$ (N/$m^3$) \\
    1 & 1 & 1 & 5 & 0.1 & 1 \\
    \hline
  \end{tabular}
  \captionof{table}{Linear and nonlinear parameters for the coupled Duffing
    system}
  \label{tab:duf_par}
\end{center}


\begin{center}
  \begin{tabular}{*{3}{c}}
    \hline
    Mode & Frequency (rad/s) & Damping ratio (\%) \\
    1 & 1.00 & 5.00 \\
    2 & 3.32 & 1.51 \\
    \hline
  \end{tabular}
  \captionof{table}{Linear natural frequencies and damping ratios for the
    coupled Duffing system}
  \label{tab:duf_eigen}
\end{center}


\begin{figure}[!ht]
  \centering
  \includegraphics[width=0.5\linewidth]{2dof_duffing/2dof_duffing}
  \caption{Schematic representation of the coupled duffing system}
  \label{fig:duf_schematic}
\end{figure}

For generating data for characterisation with WT and RFS, the system is
simulated using a logarithmic sine sweep with a sweep rate of 1 oct/min.

For identification, a single band-limited (0-15 rad/s) normally distributed
random signal with 5000 points repeated 9 times and a root mean square(rms)
value of 3N is used. This signal is called a multisine. The frequencies are
chosen to include third harmonics of the highest natural frequency.


%%% Local Variables:
%%% mode: latex
%%% TeX-master: "../report"
%%% End:
